\documentclass[12pt,a4paper,oneside]{article}
\usepackage[brazil]{babel}
\usepackage[utf8]{inputenc}
\usepackage{olymp}

\setcounter{problem}{0}

\begin{document}

\begin{problem}{Contatos}{contatos}{1}
 
José comprou um \textit{smartphone} e está com dificuldade em achar os telefones de seus contatos. O problema é que a lista de contatos está ordenada por sobrenome, algo comum em outros países. Mas ele quer seus contatos, naturalmente, ordenados pelo nome. Por exemplo, agora está assim:

\begin{center}
\begin{tabular}{l|l}
    Sobrenome & Nome \\
    \hline
    Banner & Bruce \\
    Kent & Clark \\
    Parker & Peter\\
    Wayne & Bruce
\end{tabular}
\end{center}

Mas ele quer assim:

\begin{center}
\begin{tabular}{l|l}
    Nome & Sobrenome \\
    \hline
    Bruce & Banner \\
    Bruce & Wayne \\
    Clark & Kent \\
    Peter & Parker\\
\end{tabular}
\end{center}

(note que quando há nomes iguais, eles ficam em sequência e ordenados pelo sobrenome).

Seus contatos podem até ter alguns poderes, mas eles não sabem programar, então ajude José.

%\Notes
%
%Observações, se tiver.

\InputFile

A entrada começa com uma linha contendo um número inteiro $N$, o número de contatos de José. Em seguida há $N$ linhas contendo os contatos de José no formato \texttt{sobrenome - nome - telefone}. Nome e sobrenome possuem até 15 caracteres e telefone até 10. José tem no máximo 100 contatos.

\OutputFile

Seu programa deve imprimir os contatos de José na ordem desejada, formatada conforme exemplo de saída.

% \Notes
% Preste atenção aos tipos das variáveis, pois $F(90) \approx 2^{61}$.

\Example

\begin{example}
\exmp{
4
Banner - Bruce - 999000000
Kent - Clark - 9876543210
Parker - Peter - 88008800
Wayne - Bruce - 8907866666
}{
Bruce Banner - 999000000
Bruce Wayne - 8907866666
Clark Kent - 9876543210
Peter Parker - 88008800
}%
\end{example}



%Comentários sobre o exemplo, se necessário.

\end{problem}

\end{document}
